\documentclass[12pt]{article}

\tolerance=9999
\emergencystretch=10pt
\hyphenpenalty=10000
\exhyphenpenalty=100
\usepackage[english=nohyphenation]{hyphsubst}

\usepackage{tikz} %% makes the beautiful figures in the document
\usetikzlibrary{shadows}
\usetikzlibrary{shadows.blur}

\usepackage{amsmath} %% Allows inclution of mathematical figures
\usepackage[a4paper,margin=1in]{geometry} %% Responsible for the margin and the headings of the document
\usepackage{mathptmx} %% Times new roman font
\usepackage{booktabs}
\usepackage{tabularx} %% Type of table that allows fragmentation of long sentences inside cells in tables
\usepackage{setspace}
\usepackage{multicol} %% Merging columns in tables
\usepackage{multirow} %% Merging rows in tables
\usepackage{float}
\usepackage{xcolor} %% used for coloring cells in tables
\usepackage{fancyhdr} %% Allows custom heading and footer
\usepackage{lastpage} %% Allows Referencing the last page of the document
\usepackage{calc}
\usepackage{lipsum} %% Random Words
\usepackage{ragged2e}
\usepackage{soul}
\usepackage{tcolorbox}
\usepackage{amssymb}
\usepackage{pgfplots}
\usepackage{tkz-euclide}
\usetikzlibrary{angles}
\usetikzlibrary{intersections,decorations.pathreplacing,positioning}
\usetikzlibrary{arrows.meta}
\pgfplotsset{compat=1.16}
\usepgfplotslibrary{polar}

\usepackage{hanging}

\usepackage{amsthm}

\usepgfplotslibrary{fillbetween}
\usetikzlibrary{patterns}

\usepackage{mathtools}

\usepackage{enumitem}



\newcommand*\bigcdot{\mathpalette\bigcdot@{.5}}
\newcommand*\bigcdot@[2]{\mathbin{\vcenter{\hbox{\scalebox{#2}{$\m@th#1\bullet$}}}}}


%\usepackage[scaled=1.0]{helvet}             % Load Helvetica
%\renewcommand{\familydefault}{\sfdefault}   % Set Helvetica as main font

\usepackage[T1]{fontenc}

\usepackage{gensymb}
\usepackage{orcidlink}


%%%%%%%%%%%%%%%%%%%%%%%%%%%%%%%%%%%%%%%%%%%%%%%%%%%%%%%%%%%%%%%%%%%%%%

\definecolor{h1}{HTML}{196B24}
\definecolor{h2}{HTML}{C1F0C7} %% Allows a custom color in the document using the hex code

%\pagestyle{SILPhdr} %% sets the desired page style of the entire document. Specifically, the headings and footers

\linespread{1.0}

\begin{document}
	\noindent
	SULADAY, Axelcris G. \orcidlink{https://orcid.org/0009-0007-6882-7636} \hfill PHYS102: Electromagnetic Induction and Inductance
	\rule{\textwidth}{1px}
	\begin{center}
		\hrulefill \quad \textsc{Induced emf and Magnetic Flux} \quad \hrulefill
	\end{center}
	
	\begin{hangparas}{65pt}{1}
		\textbf{Example 1:} A conducting circular loop of radius 0.250 m is placed in the \(xy\)-plane in a uniform magnetic field of 0.360 T that point in the positive \(z\)-direction, the same direction as the normal to the plane.
	\end{hangparas}
	
		\begin{proof}[Solution]
			Given the following
			\begin{align*}
				B & = 0.360~\text{T} \\
				r & = 0.250~\text{m} \\
				\theta & = 0\degree
			\end{align*}
			\begin{enumerate}
				\item[(a)] Calculate the magnetic flux through the loop.
				
					Using the formula
					\[\Phi_B = BA \cos\theta\]
					solve for the area of the loop which is
					\[A = \pi r^2 = \pi\left(0.250~\text{m}\right)^2 = \frac{1}{16}\pi~\text{m}^2\]
					then, the magnetic flux through the loop is
					\begin{align*}
						\Phi_B & = \left(0.360~\text{T}\right)\left(\frac{1}{16}\pi~\text{m}^2\right)\cos\left(0\degree\right) \\
						& = 0.07069~\text{T}\bigcdot\text{m}^2 \\
						& = 0.07069~\text{Wb}
					\end{align*}
				\item[(b)] Suppose the loop is rotated clockwise around the \(x\)-axis, so the normal direction now point at a \(45.0\degree\) angle with respect to the \(z\)-axis. Recalculate the magnetic flux through the loop.
				
					The magnetic flux through the loop is
					\begin{align*}
						\Phi_B & = \left(0.360~\text{T} \right)\left(\frac{1}{16}\pi~\text{m}^2\right)\cos\left(45\degree\right) \\
						& = 0.04998~\text{T}\bigcdot\text{m}^2 \\
						& = 0.04998~\text{Wb}
					\end{align*}
				\item[(c)] What is the change in flux due to the rotation of the loop?
				
					\[\Delta\Phi_B = 0.04998~\text{Wb} - 0.07069~\text{Wb} = -0.02071~\text{Wb}\]
			\end{enumerate}
		\end{proof}
	\newpage
		\begin{center}
			\hrulefill \quad \textsc{Faraday's Law} \quad \hrulefill
		\end{center}
		
		\begin{hangparas}{65pt}{1}
			\textbf{Example 2:} A coil with 25 turns of wire is wrapped on a frame with a square cross-section 1.80 cm on a side. Each turn has the same area, equal to that of the frame, and the total resistance of the coil is \(0.350~\ohm\). An applied uniform magnetic field is perpendicular to the plane of the coil, as shown in the figure.
		\end{hangparas}
		
		\begin{proof}[Solution]
			Given the following
			\begin{align*}
				N & = 25 \\
				A & = \left(1.80~\text{cm}\right)^2 = \left(0.0180~\text{m}\right)^2 = 3.24\times10^{-4}~\text{m}^2 \\
				R & = 0.350~\ohm
			\end{align*}
			\begin{enumerate}
				\item[(a)] If the field changes uniformly from 0.00 T to 0.500 T in 0.800 s, find the induced emf in the coil while the field is changing.
				
					using the formula
					\[\epsilon = -N \frac{\Delta\Phi_B}{\Delta t}\]
					solve the following
					\begin{align*}
						\Delta\Phi_B = B_{\perp f}A-B_{\perp i}A & = \left(0.500~\text{T}\right)\left(3.24\times10^{-4}~\text{m}^2\right)-\left(0.00~\text{T}\right)\left(3.24\times10^{-4}~\text{m}^2\right) \\
						& = 1.62\times10^{-4}~\text{Wb} \\
						\Delta t & = 0.800~\text{s}-0.00~\text{s} \\
						& = 0.800~\text{s}
					\end{align*}
					thus,
					\[\epsilon = -\left(25\right)\frac{1.62\times10^{-4}~\text{Wb}}{0.800~\text{s}} = -5.0625\times10^{-3}~\text{V}\]
					
				\item[(b)] Find the magnitude of the induced current in the coil while the field is changing.
				
					using the formula
					\[V = IR \quad \Rightarrow \quad I = \frac{V}{R} \quad \Rightarrow \quad I = \frac{|\epsilon|}{R}\]
					thus,
					\[I = \frac{|-5.0625\times10^{-3}~\text{V}|}{0.350~\ohm} = 0.014464~\text{A}\]
			\end{enumerate}
		\end{proof}
	\newpage
		\begin{center}
			\hrulefill \quad \textsc{Motional emf} \quad \hrulefill
		\end{center}
		
		\begin{hangparas}{65pt}{1}
			\textbf{Example 3:} An airplane with a wingspan of 30.0 m flies due north at a location where the downward component of the Earth's magnetic field of \(0.600\times10^{-4}\) T. There is also a component pointing due north which has a magnitude of \(0.470\times10^{-4}\) T.
		\end{hangparas}
		
		\begin{proof}[Solution]
			Given the following
			\begin{align*}
				l & = 30.0~\text{m} \\
				B & = 0.600\times10^{-4}~\text{T} \\
				v & = 2.50\times10^{2}~\text{m/s}
			\end{align*}
			\begin{enumerate}
				\item[(a)] Find the difference in potential between the wingtips when the speed of the plane is \(2.50\times10^2\) m/s.
				
					using the formula
					\[V_{ab} = vBL\]
					then, the difference in potential is
					\begin{align*}
						V_{ab} & = \left(2.50\times10^2~\text{m/s}\right)\left(0.600\times10^{-4}~\text{T}\right)\left(30.0~\text{m}\right) \\
						& = 0.450~\text{V}
					\end{align*}
			\end{enumerate}
		\end{proof}
	\newpage
	
	\begin{center}
		\hrulefill \quad \textsc{Self-Inductance} \quad \hrulefill
	\end{center}
	
	\begin{hangparas}{65pt}{1}
		\textbf{Example 4:} Calculate the inductance of solenoid containing 300 turns if the lenth of the solenoid is 25.0 cm and its cross-sectional area is \(4.00\times10^{-4}~\text{m}^2\). (b) Calculate the self-induced emf in the solenoid described in (a) if the current in the solenoid decreases at the rate of 50.0 A/s.
	\end{hangparas}
	
	\begin{proof}[Solution]
	
	\end{proof}
	
\end{document}